\documentclass[11pt,a4paper]{article}
\usepackage{pkpaper}
% generated by scripts/measure.py; do not edit
\newcommand{\pkItems}{44}
\newcommand{\pkItemsStd}{16}
\newcommand{\pkItemsAuthor}{28}
\newcommand{\pkItemsRewrite}{38}
\newcommand{\pkItemsConv}{4}
\newcommand{\pkItemsGap}{2}
\newcommand{\pkProbes}{420}
\newcommand{\pkProbesUnique}{408}
\newcommand{\pkRules}{32}
\newcommand{\pkRulesR}{21}
\newcommand{\pkRulesRK}{10}
\newcommand{\pkRulesK}{1}
\newcommand{\pkRulesRegular}{16}
\newcommand{\pkRulesSyntactic}{11}
\newcommand{\pkRulesSemantic}{5}
\newcommand{\pkRulesLinted}{19}
\newcommand{\pkLintChecks}{21}
\newcommand{\pkLintDetected}{26}
\newcommand{\pkLintUndetected}{18}
\newcommand{\pkLintPkLongest}{15}
\newcommand{\pkCatalanNineteen}{1\,767\,263\,190}
\newcommand{\pkCatalanTwentyFour}{1\,289\,904\,147\,324}
\newcommand{\pkAttachTen}{58\,786}
\newcommand{\pkAttachEight}{4862}
\newcommand{\pkBracketEleven}{4066}
\newcommand{\pkBracketElevenAll}{16\,796}
\newcommand{\pkScopeSeven}{1066}
\newcommand{\pkScopeSevenAll}{5040}
\newcommand{\pkScopeFour}{14}
\newcommand{\pkScopeFourAll}{24}
\newcommand{\pkOpenList}{256}
\newcommand{\pkParseExpNat}{2,53}
\newcommand{\pkParseExpPk}{1,77}
\newcommand{\pkParseWordsMax}{83}
\newcommand{\pkParseOpsNatMax}{15\,345}
\newcommand{\pkParseOpsPkMax}{1976}
\newcommand{\pkParseRatioMax}{7,8}
\newcommand{\pkParseOpsNatTwenty}{496}
\newcommand{\pkParseOpsPkTwenty}{178}
\newcommand{\pkParseTreesTwenty}{16\,796}
\newcommand{\pkParseRatioTwentyFive}{3,1}
\newcommand{\pkParseTreesTwentyFive}{208\,012}
\newcommand{\pkItemsPkSentences}{88}
\newcommand{\pkTokItems}{40}
\newcommand{\pkTokMedMin}{36}
\newcommand{\pkTokMedMax}{39}
\newcommand{\pkProbesGrayUnique}{136}
\newcommand{\pkClaudeCalls}{272}
\newcommand{\pkClaudeTokMedian}{687}
\newcommand{\pkPrecN}{60}
\newcommand{\pkPrecValid}{57}
\newcommand{\pkPrec}{95}
\newcommand{\pkPrecLo}{86}
\newcommand{\pkPrecHi}{98}
\newcommand{\pkRedInstances}{400}
\newcommand{\pkRedSat}{398}
\newcommand{\pkRedUnsat}{2}
\newcommand{\pkTokUnit}{o200k}
\newcommand{\pkTokTokenizers}{4}
\newcommand{\pkTokFailed}{0}
\newcommand{\pkTokMedian}{37}
\newcommand{\pkTokQone}{13}
\newcommand{\pkTokQthree}{63}
\newcommand{\pkTokMean}{49}
\newcommand{\pkTokMin}{−4}
\newcommand{\pkTokMax}{286}
\newcommand{\pkTokCiLo}{23}
\newcommand{\pkTokCiHi}{53}
\newcommand{\pkTokShareLonger}{90}
\newcommand{\pkTokBreakEven}{27}
\newcommand{\pkTokBreakEvenMean}{33}
\newcommand{\pkWordMedian}{38}
\newcommand{\pkWordMean}{52}
\newcommand{\pkCharMedian}{39}
\newcommand{\pkTokDeltaMedian}{0,37}
\newcommand{\pkTokDeltaMean}{0,49}
\newcommand{\pkLegendTokMax}{37}
\newcommand{\pkProbeModels}{6}
\newcommand{\pkProbeModelsGraded}{4}
\newcommand{\pkProbeLocal}{4}
\newcommand{\pkProbeClaude}{2}
\newcommand{\pkProbeQualified}{5}
\newcommand{\pkProbeFailed}{1}
\newcommand{\pkAccPkMin}{54}
\newcommand{\pkAccPkMax}{99}
\newcommand{\pkAgreeItems}{44}
\newcommand{\pkNatSplit}{73}
\newcommand{\pkPkAllCorrect}{27}
\newcommand{\pkNatSplitCount}{32}
\newcommand{\pkPkAllCorrectCount}{12}
\newcommand{\pkErrPk}{17}
\newcommand{\pkBreakEvenStar}{40}
\newcommand{\pkAnchorMin}{65}
\newcommand{\pkAnchorMax}{92}
\newcommand{\pkAccPkQMin}{70}
\newcommand{\pkAccPkQMax}{99}
\newcommand{\pkSteerQMin}{0,19}
\newcommand{\pkSteerQMax}{0,97}
\newcommand{\pkUndecidedQMin}{0}
\newcommand{\pkUndecidedQMax}{86}
\newcommand{\pkConvPlainMin}{50}
\newcommand{\pkConvPlainMax}{75}
\newcommand{\pkConvLegendMin}{50}
\newcommand{\pkConvLegendMax}{100}
\newcommand{\pkConvOrderAll}{20}
\newcommand{\pkConvOrderPlain}{14}
\newcommand{\pkConvOrderLegend}{15}
\newcommand{\pkConvOrAll}{20}
\newcommand{\pkConvOrPlain}{10}
\newcommand{\pkConvOrLegend}{15}
\newcommand{\pkProbeQualifiedConv}{5}
\newcommand{\pkGapOneReadingOne}{1}
\newcommand{\pkGapOnePkOk}{2}
\newcommand{\pkGapTwoReadingOne}{5}
\newcommand{\pkGapTwoPkOk}{2}
\newcommand{\pkHNatMin}{0,00}
\newcommand{\pkHNatMax}{0,91}
\newcommand{\pkHPkMin}{0,00}
\newcommand{\pkHPkMax}{0,89}
\newcommand{\pkSignWins}{0}
\newcommand{\pkProbeLocalQualified}{3}
\newcommand{\pkClaudeItems}{44}
\newcommand{\pkClaudeSplitCount}{5}
\newcommand{\pkClaudeAgreeCount}{30}
\newcommand{\pkClaudeUnsureCount}{9}
\newcommand{\pkClaudePkBothCount}{34}
\newcommand{\pkAccPkClaudeMin}{91}
\newcommand{\pkAccPkClaudeMax}{99}
\newcommand{\pkSteerClaudeMin}{0,84}
\newcommand{\pkSteerClaudeMax}{0,97}
\newcommand{\pkLocalParamMin}{1,6}
\newcommand{\pkLocalParamMax}{4,8}
\newcommand{\pkAccPkLocalMin}{54}
\newcommand{\pkAccPkLocalMax}{82}
\newcommand{\pkSteerLocalMin}{0,12}
\newcommand{\pkSteerLocalMax}{0,50}
\newcommand{\pkHNatLocalMin}{0,00}
\newcommand{\pkHNatLocalMax}{0,91}
\newcommand{\pkHPkLocalMin}{0,00}
\newcommand{\pkHPkLocalMax}{0,89}
\newcommand{\pkSignWinsLocal}{0}
\newcommand{\pkUnstableNat}{10}
\newcommand{\pkUnstablePk}{1}
\newcommand{\pkCorpusStdCfive}{2}
\newcommand{\pkCorpusActs}{5}
\newcommand{\pkCorpusPk}{3}
\newcommand{\pkCorpusActWords}{219\,801}
\newcommand{\pkCorpusActMin}{43}
\newcommand{\pkCorpusActMax}{71}
\newcommand{\pkCorpusPkMin}{0,0}
\newcommand{\pkCorpusPkMax}{3,5}
\newcommand{\pkCorpusOverMin}{12}
\newcommand{\pkCorpusOverMax}{45}


\title{\textbf{Sprawozdanie z pomiarów}\\[4pt]
\large do pracy \pkq{Polszczyzna kontrolowana jako metoda redukcji\\ złożoności obliczeniowej w AI}}
\author{Maciej Szulia\\[2pt]\normalsize\texttt{maciej.szulia@gmail.com}}
\date{2026-11-12}

\begin{document}
\maketitle

\section{Zakres}

Sprawozdanie opisuje 5 pomiarów pracy: warunki, wyniki, odstępstwa od planu i sposób
odtworzenia. Autor napisał sprawozdanie w PK. Legenda: \pkq{lub} znaczy co najmniej 1
z członów; \pkq{albo} znaczy dokładnie 1 z członów. Znacznik w nawiasie kwadratowym podaje
źródło faktu. Każda liczba w tabelach pochodzi z plików katalogu \texttt{results}
\pkq{[obliczenie: scripts/measure.py]}.

\section{Środowisko}

\begin{center}
\begin{tabular}{ll}
\toprule
Element & Wartość\\
\midrule
Data pomiarów & 2026-11-12\\
Komputer & procesor Intel Core i7-7700HQ, 8 wątków, 31 GB pamięci\\
System & Linux 7.2.9, dystrybucja Artix\\
Python & 3.13.14\\
Program \texttt{ollama} & 0.34.1, obliczenia na procesorze\\
Polecenie \texttt{claude} & 2.1.322\\
Przedmiot badania & standard PK 0.1, wersja \texttt{15503a4}\\
\bottomrule
\end{tabular}
\end{center}
Źródło danych tabeli: \pkq{[obserwacja: 2026-11-12, 18:31]}.

\section{Wyniki zbiorcze}

\begin{center}
\small
\begin{tabular}{@{}l>{\raggedright\arraybackslash}p{0.31\textwidth}>{\raggedright\arraybackslash}p{0.35\textwidth}l@{}}
\toprule
Pomiar & Wielkość & Wynik & Plik wyników\\
\midrule
P1 & liczba odczytań zdania z 8 frazami przyimkowymi & \pkAttachEight{} odczytania; zdanie PK: 1 & \texttt{p1\_counts.json}\\
P2 & wykładnik wzrostu liczby operacji parsera & \pkParseExpNat{} bez reguł PK; \pkParseExpPk{} z regułą B6 & \texttt{p2\_parser.json}\\
P3 & narzut długości PK w tokenach \pkTokUnit & mediana \pkTokMedian{} \%; kwartyle \pkTokQone{} \% i \pkTokQthree{} \% & \texttt{p3\_tokens.json}\\
P4 & trafność odczytania zdań PK przez model & od \pkAccPkQMin{} \% do \pkAccPkQMax{} \% włącznie (modele po kontroli kotwic) & \texttt{p4\_summary.json}\\
P4 & rozbieżność modeli na zdaniu zwykłym & \pkNatSplitCount{} z \pkAgreeItems{} pozycji & \texttt{p4\_summary.json}\\
P5 & trafienia sprawdzarki na 1000 słów & akty prawne: od \pkCorpusActMin{} do \pkCorpusActMax{} włącznie; teksty PK: od \pkCorpusPkMin{} do \pkCorpusPkMax{} włącznie & \texttt{p5\_corpus.json}\\
P5 & precyzja sprawdzarki & \pkPrecValid{} z \pkPrecN{} trafień, czyli \pkPrec{} \% & \texttt{summary.json}\\
\bottomrule
\end{tabular}
\end{center}

Kontrole poprawności: test redukcji z problemu 3-SAT przeszedł \pkRedInstances{} z
\pkRedInstances{} formuł. Sprawdzarka wykrywa \pkLintDetected{} z \pkItems{} zdań zwykłych
zbioru testowego i nie zgłasza żadnego z \pkItemsPkSentences{} zdań PK.

\section{Pomiar P4: biegi i wyniki}

Zbiór ma \pkItems{} pozycji: \pkItemsRewrite{} przeredagowań, \pkItemsConv{} konwencje
i \pkItemsGap{} luki. Pełna sonda ma \pkProbesUnique{} różnych zapytań na model. Tor Claude
obejmuje \pkProbesGrayUnique{} zapytań strefy szarej, każde w 2 próbach.

\begin{table}[htp]
\centering
\caption{Biegi sondy. Czas obejmuje nowe wywołania 1 biegu; bieg wznowiony ma osobny wiersz.}
\label{tab:runs}
\small
\begin{tabular}{llrr}
\toprule
Tor & Model & nowe zapytania & czas (min)\\
\midrule
\wiersze{data/table_runs.tex}
\bottomrule
\end{tabular}
\end{table}

\begin{table}[htp]
\centering
\caption{Wynik sondy. L: tor lokalny, C: tor Claude. $n$: liczba pozycji typu
przeredagowanie. $H$: mediana niepewności w bitach. Znaki: liczba pozycji z niepewnością
wyższą dla zdania zwykłego i liczba pozycji z niepewnością wyższą dla zdań PK; $p$: test
znaków. Kotwice: trafność w procentach; próg kwalifikacji modelu wynosi 80 \%.}
\label{tab:probe}
\small
\setlength{\tabcolsep}{4.5pt}
\begin{tabular}{@{}lrrrrrrrrr@{}}
\toprule
 & & \multicolumn{2}{c}{$H$ (bity)} & \multicolumn{2}{c}{test znaków} & & trafność & \multicolumn{2}{c}{kotwice (\%)}\\
Model & $n$ & zwykłe & PK & znaki & $p$ & sterow. & PK (\%) & zwykłe & PK\\
\midrule
\wiersze{data/table_probe.tex}
\bottomrule
\end{tabular}
\end{table}

\begin{table}[htp]
\centering
\caption{Pozycje typu konwencja (\pkItemsConv{} pozycje, reguły B2 i C1): trafność
w procentach i sterowalność bez legendy oraz z legendą.}
\label{tab:conv}
\small
\begin{tabular}{lrrrr}
\toprule
Model & trafność & trafność & sterowalność & sterowalność\\
 & bez legendy & z legendą & bez legendy & z legendą\\
\midrule
\wiersze{data/table_conv.tex}
\bottomrule
\end{tabular}
\end{table}

Kontrolę kotwic przeszło \pkProbeQualified{} z \pkProbeModels{} modeli; modele toru Claude
nie mają kotwic i wchodzą do zestawienia na podstawie trafności PK. W \pkNatSplitCount{}
z \pkAgreeItems{} pozycji modele po kontroli wybierają różne odczytania zdania zwykłego.
W \pkPkAllCorrectCount{} z \pkAgreeItems{} pozycji każdy model po kontroli odczytuje oba
zdania PK zgodnie z kluczem.

\section{Pomiary P3 i P5}

\begin{table}[htp]
\centering
\caption{P3: narzut długości PK w procentach długości zdania zwykłego; \pkTokItems{}
pozycji.}
\small
\begin{tabular}{lrrrrrr}
\toprule
Miara & mediana & kwartyl 1 & kwartyl 3 & minimum & maksimum & dłuższe (\%)\\
\midrule
\wiersze{data/table_tokens.tex}
\bottomrule
\end{tabular}
\end{table}

\begin{table}[htp]
\centering
\caption{P5: wynik sprawdzarki na korpusie. Mediana: mediana długości zdania w słowach.
Kolumna 4: udział zdań dłuższych niż 25 słów. Trafienia: liczba na 1000 słów. Kolumna końcowa:
wystąpienia słowa \pkq{się} na 1000 słów, kandydat do oceny.}
\small
\begin{tabular}{@{}lrrrrr@{}}
\toprule
Tekst & słowa & mediana & ponad 25 (\%) & trafienia & \pkq{się}\\
\midrule
\wiersze{data/table_corpus.tex}
\bottomrule
\end{tabular}
\end{table}

Tokenizery modeli Bielik wymagają zgody licencyjnej w serwisie Hugging Face; żądanie pliku
zwróciło błąd 401 \pkq{[obserwacja: 2026-11-12, 17:52]}. Pomiar P3 obejmuje więc
\pkTokTokenizers{} tokenizery.

\section{Odstępstwa od planu i problemy}

\begin{enumerate}[leftmargin=1.8em,itemsep=2pt]
  \item \textbf{Tor OpenRouter nie powstał.} Plan przewidywał 3 tory pomiaru P4.
  Zabezpieczenie komputera autora zablokowało odczyt klucza dostępu
  \pkq{[obserwacja: 2026-11-12, 15:15]}. Pomiar ma 2 tory: lokalny i Claude.
  \item \textbf{Zbiór ma \pkItems{} pozycji.} Plan przewidywał 40 pozycji. Dodatkowe
  pozycje to 2 luki standardu (G1, G2) i 2 pozycje o słowach kluczowych (K1, K2).
  \item \textbf{Kontrola kotwic.} Kontroli kotwic nie przeszło \pkProbeFailed{}
  z \pkProbeLocal{} modeli lokalnych. Wynik takiego modelu stoi w tabeli \ref{tab:probe}
  i nie wchodzi do miar zgodności.
  \item \textbf{Tor Claude nie ma kotwic.} Budżet toru wynosił co najwyżej 600 wywołań,
  więc tor obejmuje tylko sondy strefy szarej w 2 próbach.
  \item \textbf{Przerwa toru Claude.} O godz. 18:06 polecenie \texttt{claude} zniknęło na
  czas aktualizacji pakietu i bieg zakończył się błędem. Autor dodał do skryptu ponawianie
  i wznowił bieg; skrypt zachował odpowiedzi sprzed przerwy
  \pkq{[obserwacja: 2026-11-12, 18:07]}. Bieg przerwany wykonał 34 zapytania i nie ma
  wiersza w tabeli \ref{tab:runs}.
  \item \textbf{Odczyt werdyktu w torze Claude.} W 4 sondach model napisał rachunek przed
  werdyktem, a skrypt odczytywał werdykt z początku odpowiedzi. Autor zmienił regułę:
  werdyktem jest końcowe słowo odpowiedzi. Autor powtórzył te 4 sondy (8 wywołań)
  i przeliczył 1 wynik z zapisanej odpowiedzi. Tor Claude zużył razem 552 wywołania.
  \item \textbf{Tokenizery modeli Bielik.} Serwis Hugging Face zwrócił błąd 401, więc
  pomiar P3 pomija 2 tokenizery.
  \item \textbf{Pobieranie modeli.} Łącze dało prędkość 1,3 MB/s \pkq{[obserwacja:
  2026-11-12, 17:41]}. Program \texttt{ollama} zatrzymał pobieranie 1 raz; skrypt ponawia
  pobieranie co 10 minut. Modele zajmują miejsce na partycji \texttt{/mnt/windows}, która
  przed pobraniem miała 18 GB wolnego miejsca.
  \item \textbf{Stan modeli lokalnych.} Pomiar objął \pkProbeLocal{} modele lokalne z 4 planowanych: \texttt{llama3.2:3b},
\texttt{gemma3:4b}, Bielik 1.5B i Bielik 4.5B. Pobranie 4 modeli zajęło 2 h 42 min
\pkq{[obserwacja: 2026-11-12, od 17:27 do 20:09]}.

  \item \textbf{Parser bez poprawki Leo.} Gramatyka PK w pomiarze P2 ma wykładnik
  \pkParseExpPk{}, a nie 1. Czas liniowy wymaga parsera z poprawką Leo.
  \item \textbf{Sprawdzarka jest przybliżona.} Sprawdzarka zgłasza także fałszywe
  trafienia, na przykład przymiotnik w nazwie \pkq{małe litery}. Precyzja na próbie
  \pkPrecN{} trafień wynosi \pkPrec{} \%. Czułość jest nieznana.
  \item \textbf{Podział aktów prawnych na zdania.} Akt składa się z jednostek
  redakcyjnych. Skrypt usuwa numery jednostek i dzieli tekst regułą interpunkcyjną, więc
  mediana długości zdania jest przybliżeniem.
\end{enumerate}


\section{Odtworzenie}

\begin{itemize}[leftmargin=1.6em,itemsep=1pt]
  \item \texttt{make all}: testy, pomiary P1, P2, P3 i P5, agregaty pomiaru P4, rysunki
  i pliki PDF. Pomiary P3 i P5 pobierają dane z sieci.
  \item \texttt{make MEASURE\_FLAGS=-{}-offline all}: to samo bez sieci, z plików
  \texttt{results/p3\_tokens.json} i \texttt{results/p5\_corpus.json}.
  \item \texttt{python scripts/probe.py ollama MODEL}: pełna sonda na modelu lokalnym;
  skrypt pomija zapytania, które plik wyników już zawiera.
  \item \texttt{python scripts/probe.py claude MODEL -{}-subset gray -{}-samples 2}:\\ sonda toru Claude.
\end{itemize}
Surowe odpowiedzi modeli zawierają pliki \texttt{results/probe\_*.jsonl}, a dzienniki
biegów zawiera katalog \texttt{results/logs}.

\end{document}
